% Einheit 2 von 4 – Addieren und Subtrahieren (bank/rationale-zahlen/e2.jsonl, zusammenbau v0.1)
%% TODO Zweigzeile Teil 1: was hier gelernt wird (Satz in Schülersprache aus dem Einheitstitel) – Einheit 2
\einheitenkopf[e2]{Einheit 2 von 4 · Addieren und Subtrahieren}
\zweigzeile{ab Kl. 6, je nach Buch bis Kl. 7 (am Gymnasium ab Kl. 7) · P10}
\setcounter{aufgabe}{12}

%% TODO Titel: Ich-kann-Satz fehlt in der Bank (Platzhalter: Kettenname) – Nr. 13
%% TODO Anweisung: ein Satz über der ersten Teilaufgabe fehlt in der Bank – Nr. 13
\begin{aufgabe}{Vorzeichen oder Rechenzeichen?}
\begin{teile}
\teil $-4 - 6$ – jedes Minus: Vorzeichen (V) oder Rechenzeichen (R)? \leerfeld
\end{teile}
\end{aufgabe}

%% TODO Titel: Ich-kann-Satz fehlt in der Bank (Platzhalter: Kettenname) – Nr. 14
%% TODO Anweisung: ein Satz über der ersten Teilaufgabe fehlt in der Bank – Nr. 14
\begin{aufgabe}{Zeichen zusammenfassen}
\begin{teile}
\teil $6 - (-5)$ – ohne Klammern? \leerfeld
\end{teile}
\end{aufgabe}

%% TODO Titel: Ich-kann-Satz fehlt in der Bank (Platzhalter: Kettenname) – Nr. 15
%% TODO Anweisung: ein Satz über der ersten Teilaufgabe fehlt in der Bank – Nr. 15
\begin{aufgabe}{Addieren/Subtrahieren}
\swa{Zeichne den Pfeil: Start bei $-4$, $6$ nach rechts.}{\zahlenstrahl[xmin=-8,xmax=8,xstep=1]{}}
%% TODO Darstellung für schwach fehlt in der Bank (rationale-zahlen-e2-k3-s1-v1; Abschnitt „Für schwache Schüler“)
\swz{Berechne: $-5 + 8$}{}
%% TODO Darstellung für schwach fehlt in der Bank (rationale-zahlen-e2-k3-s1-v2; Abschnitt „Für schwache Schüler“)
\swz{Berechne: $-3 + 2$}{}
%% TODO Darstellung für schwach fehlt in der Bank (rationale-zahlen-e2-k3-s1-v3; Abschnitt „Für schwache Schüler“)
\swz{Berechne: $-4 - 3$}{}
%% TODO Darstellung für schwach fehlt in der Bank (rationale-zahlen-e2-k3-s1-v4; Abschnitt „Für schwache Schüler“)
\swz{Berechne: $-9 + 4$}{}
%% TODO Darstellung für schwach fehlt in der Bank (rationale-zahlen-e2-k3-s1-v5; Abschnitt „Für schwache Schüler“)
\swz{Morgens $-7$ °C, dann $5$ Grad wärmer – wie warm?}{}
\swfrage{Was bleibt gleich, was ändert sich?}
%% TODO Darstellung für schwach fehlt in der Bank (rationale-zahlen-e2-k3-s2-v1; Abschnitt „Für schwache Schüler“)
\swz{Berechne: $7 + (-3)$}{}
%% TODO Darstellung für schwach fehlt in der Bank (rationale-zahlen-e2-k3-s3-v1; Abschnitt „Für schwache Schüler“)
\swz{Berechne: $6 - (-7)$}{}
%% TODO Darstellung für schwach fehlt in der Bank (rationale-zahlen-e2-k3-s4-v1; Abschnitt „Für schwache Schüler“)
\swz{Berechne: $3 - (+8) + (-4)$}{}
%% TODO Darstellung für schwach fehlt in der Bank (rationale-zahlen-e2-k3-s5-v1; Abschnitt „Für schwache Schüler“)
\swz{$-6$ und $5$ – wie weit auseinander?}{}
%% TODO Darstellung für schwach fehlt in der Bank (rationale-zahlen-e2-k3-s6-v1; Abschnitt „Für schwache Schüler“)
\swz{Berechne: $-2{,}5 + 4{,}1$}{}
%% TODO Darstellung für schwach fehlt in der Bank (rationale-zahlen-e2-k3-s7-v1; Abschnitt „Für schwache Schüler“)
\swz{Berechne: $-3 + 8 - 15 + 4$}{}
\swz[3]{Berechne die Klammer von $(a + b) : c$ für $a = 3$ und $b = -17$.}{}
\end{aufgabe}

%% TODO Titel: Ich-kann-Satz fehlt in der Bank (Platzhalter: Kettenname) – Nr. 16
%% TODO Anweisung: ein Satz über der ersten Teilaufgabe fehlt in der Bank – Nr. 16
\begin{aufgabe}{Beträge: gleiche Vorzeichen addieren, verschiedene subtrahieren}
%% TODO Darstellung für schwach fehlt in der Bank (rationale-zahlen-e2-k4-s1-v1; Abschnitt „Für schwache Schüler“)
\swz{$-17 + (-8)$ – gleiche oder verschiedene Vorzeichen? Rechne mit den Beträgen.}{}
\end{aufgabe}

%% TODO Titel: Ich-kann-Satz fehlt in der Bank (Platzhalter: Kettenname) – Nr. 17
%% TODO Anweisung: ein Satz über der ersten Teilaufgabe fehlt in der Bank – Nr. 17
\begin{aufgabe}{Dezimalzahlen und Brüche}
%% TODO Darstellung für schwach fehlt in der Bank (rationale-zahlen-e2-k5-s1-v1; Abschnitt „Für schwache Schüler“)
\swz{Berechne: $-\frac{3}{4} + \frac{1}{4}$}{}
\end{aufgabe}

%% TODO Titel: Ich-kann-Satz fehlt in der Bank (Platzhalter: Kettenname) – Nr. 18
%% TODO Anweisung: ein Satz über der ersten Teilaufgabe fehlt in der Bank – Nr. 18
\begin{aufgabe}{Umkehrung: Rechnung zu Ergebnis finden}
%% TODO Darstellung für schwach fehlt in der Bank (rationale-zahlen-e2-k6-s1-v1; Abschnitt „Für schwache Schüler“)
\swz{$-7 + \square = 2$ – welche Zahl fehlt?}{}
\end{aufgabe}

%% TODO Titel: Ich-kann-Satz fehlt in der Bank (Platzhalter: Kettenname) – Nr. 19
%% TODO Anweisung: ein Satz über der ersten Teilaufgabe fehlt in der Bank – Nr. 19
\begin{aufgabe}{Addieren/Subtrahieren – Fehler finden, Begründen, Darstellungswechsel, Anwendung}
\swz[3]{Ole rechnet: $-4 - 7 = -3$. Finde den Fehler.}{}
\swfrage{Setze fort: $4 - 2 = 2$, $4 - 1 = 3$, $4 - 0 = 4$. Begründe damit, warum $4 - (-3)$ dasselbe ist wie $4 + 3$.}
\swz{Ein Pfeil startet bei $-3$ und endet bei $4$. Welche Rechnung?}{\zahlenstrahl[xmin=-5,xmax=5,xstep=1]{-3/Start, 4/Ende}}
\swz[3]{Ein Aufzug steht in der Tiefgarage auf Ebene $-2$ und fährt $7$ Stockwerke nach oben. Auf welcher Ebene hält er?}{}
\end{aufgabe}

\uebersichtskasten{Zeichen zusammenfassen: gleiche Zeichen werden plus, verschiedene werden minus. \\ 9 − (−4) = 9 + 4 = 13 \quad 9 + (−4) = 9 − 4 = 5 \\ Dann an der Zahlengeraden: plus nach rechts, minus nach links. \\ −6 + 10 = 4 \quad −6 − 10 = −16}
